\chapter*{ABSTRACT}
\addcontentsline{toc}{chapter}{ABSTRACT}

University leadership at every level needs timely, trustworthy visibility into research performance, yet research data is scattered across indexing databases and internal records, and reports are often compiled manually. This project presents the Research Intelligence Agent (RIA), a multi-level research analytics and decision-support platform built for KLE Technological University. RIA stores 6,516 publication records and 739 faculty profiles across 18 institutes and 36 departments in a PostgreSQL database managed through Supabase, and aggregates metrics upward through the organizational hierarchy.

The platform defines five role levels: L1 Vice Chancellor, L2 Executive Dean, L3 Faculty Dean, L4 Head of Department, and L5 Faculty. A role-specific AI persona defines each level's goals, communication tone, permitted tool scope, and response boundaries. At login, the authenticated user's role selects the matching persona, scope filter, and tool set from five persona modules (\texttt{persona\_l1} through \texttt{persona\_l5}). Access scope is enforced server-side through \texttt{scope\_guard.py}: a two-layer system that (1) filters which of the 17 analytics tools the LLM is shown at all, and (2) rewrites or injects trusted arguments---such as \texttt{ria\_user\_id} and \texttt{scope\_departments}---into every tool call before execution, making scope enforcement unbypassable by prompt injection.

The backend is implemented in FastAPI with a Groq LLM (accessed via OpenRouter) for the agentic tool-calling loop, and an LLM-free deterministic intent router that handles common queries in under 15\,ms at zero token cost. The frontend is built with React 19 and Vite 7, using Tailwind CSS and Recharts for role-tailored dashboards at five dedicated routes. A 53-case test suite covering frontend, backend/API, agentic AI, and end-to-end flows was executed and all cases passed, with university aggregate queries returning 3,395 publications matching direct SQL ground truth.

This report describes the problem, requirements, architecture, data model, import pipelines, persona design, and validated test results. The project makes research reporting continuous, explainable, and less dependent on any single individual while preserving privacy across organizational levels through server-side scope enforcement rather than prompt wording alone.

\noindent\textbf{Keywords:} Research analytics, role-based access control, agentic AI, persona design, tool calling, Supabase, PostgreSQL, FastAPI, Groq LLM, Tailwind CSS, Vite, data pipeline, publication analytics, scope enforcement.
